


\addcontentsline{toc}{section}{Приложение С. Код заполнения таблиц второго и третьего уровней}
\section*{Приложение С. Код заполнения таблиц второго и третьего уровней}
\begin{Verbatim}[numbers=left, frame=single, breaklines=true, breakanywhere=true]

import random
from faker import Faker
import psycopg2
from datetime import datetime, date, timedelta, time as datetime_time

fake = Faker('ru_RU')

conn = psycopg2.connect(
    host="localhost",
    user="julia7",
    password="St080806sql",
    database="cows_db"
)
cursor = conn.cursor()

total_records = {
    'Breed': 0,
    'Day_Times': 0,
    'Season': 0,
    'Lactation_Phase': 0,
    'Feed_Type': 0,
    'Livestock_specialist': 0,
    'Cow': 0,
    'Feeds': 0,
    'Diet': 0,
    'Feeding_plan': 0,
    'Nutrition': 0,
    'Lactation': 0
}

def create_cows(count=6574):
    # Получаем список всех ID пород из таблицы Breed
    cursor.execute("SELECT id_Breed FROM Breed")
    breed_ids = [row[0] for row in cursor.fetchall()]
    
    if not breed_ids:
        print("Нет пород в таблице Breed! Сначала заполните Breed.")
        return
    
    num_breeds = len(breed_ids)  # будет 10
    cows_per_breed = count // num_breeds  
    remainder = count % num_breeds  
    
    for idx, breed_id in enumerate(breed_ids):
        current_count = cows_per_breed + (1 if idx < remainder else 0)
        
        for _ in range(current_count):
            weight = random.randint(350, 850)
            age = random.randint(1, 15)
            
            cursor.execute(
                "INSERT INTO Cow (Weight, Age, id_Breed) VALUES (%s, %s, %s)",
                (weight, age, breed_id)
            )
            total_records['Cow'] += 1
    
    conn.commit()
    print(f"Добавлено {total_records['Cow']} коров")

def create_feeds(count=40):
    # Получаем все id типов кормов
    cursor.execute("SELECT id_Feed_type FROM Feed_Type")
    feed_type_ids = [row[0] for row in cursor.fetchall()]

    if not feed_type_ids:
        print("Нет данных в Feed_Type! Сначала заполните словарь.")
        return

    feed_names = [
        "Сено люцерны", "Сено луговое", "Сенаж клеверный", "Силос кукурузный",
        "Солома пшеничная", "Солома ячменная", "Свекла кормовая",
        "Морковь кормовая", "Зерно кукурузы", "Зерно ячменя",
        "Комбикорм ПК-60", "Отруби пшеничные", "Жмых подсолнечный",
        "Шрот соевый", "Патока свекловичная", "Премикс витаминный",
        "Соль поваренная", "Мел кормовой", "Мясокостная мука",
        "Зеленая трава", "ОАК концентрат"
    ]

    physical_chars = [
        "Сыпучий, сухой", "Влажный, плотный", "Грубая структура",
        "Мелкодисперсный", "Волокнистый", "Гранулированный"
    ]

    chemical_props = [
        "Высокое содержание белка",
        "Богат углеводами",
        "Повышенное содержание клетчатки",
        "Минерально-витаминный состав",
        "Высокая энергетическая ценность",
        "Низкое содержание жира"
    ]

    for _ in range(count):
        name = random.choice(feed_names) + f" {random.randint(1,100)}"

        cursor.execute(
            """
            INSERT INTO Feeds (Name, Physical_characteristics, Chemical_properties, id_Feed_type)
            VALUES (%s, %s, %s, %s)
            """,
            (
                name,
                random.choice(physical_chars),
                random.choice(chemical_props),
                random.choice(feed_type_ids)
            )
        )

        total_records['Feeds'] += 1

    conn.commit()
    print(f"Добавлено {total_records['Feeds']} кормов")

def create_diets():
    # Получаем id зоотехников
    cursor.execute("SELECT id_Livestock_specialist FROM Livestock_specialist")
    specialist_ids = [row[0] for row in cursor.fetchall()] # получение списка зоотехников

    if not specialist_ids:
        print("Нет зоотехников!")
        return

    # Получаем id сезонов
    cursor.execute("SELECT id_Season FROM Season")
    season_ids = [row[0] for row in cursor.fetchall()] # получение списка сезонов

    # Получаем id фаз лактации
    cursor.execute("SELECT id_Lactation_phase FROM Lactation_Phase")
    phase_ids = [row[0] for row in cursor.fetchall()] # получение списка фаз

    if not season_ids or not phase_ids:
        print("Не заполнены словари Season или Lactation_Phase!")
        return

    for specialist_id in specialist_ids:
        diets_count = random.randint(1, 10)  # от 1 до 10 рационов для одного специалиста

        for _ in range(diets_count):
            cursor.execute(
                """
                INSERT INTO Diet (id_Season, id_Lactation_phase, id_Livestock_specialist)
                VALUES (%s, %s, %s)
                """,
                (
                    random.choice(season_ids),
                    random.choice(phase_ids),
                    specialist_id
                )
            )

            total_records['Diet'] += 1

    conn.commit()
    print(f"Добавлено {total_records['Diet']} рационов")

def create_lactations():
    # Получаем коров
    cursor.execute("SELECT id_Cow FROM Cow")
    cow_ids = [row[0] for row in cursor.fetchall()]

    if not cow_ids:
        print("Нет коров!")
        return

    # Получаем фазы лактации (важно: порядок!)
    cursor.execute("""
        SELECT id_Lactation_phase 
        FROM Lactation_Phase
        ORDER BY id_Lactation_phase
    """)
    phase_ids = [row[0] for row in cursor.fetchall()]

    if len(phase_ids) != 5:
        print("Ожидается 5 фаз лактации!")
        return

    for cow_id in cow_ids:
        lactation_count = random.randint(1, 7)

        # начальная дата (в прошлом)
        current_date = fake.date_between(start_date='-5y', end_date='-1y')

        for lact_num in range(1, lactation_count + 1):

            # длительности фаз (в днях, примерно реалистично)
            phase_durations = [
                random.randint(10, 30),   # новотельный период
                random.randint(80, 150),  # середина
                random.randint(30, 60),   # заключительная
                random.randint(20, 40),   # ранний сухостой
                random.randint(20, 40)    # поздний сухостой
            ]

            for phase_id, duration in zip(phase_ids, phase_durations):
                start_date = current_date
                end_date = start_date + timedelta(days=duration)

                cursor.execute(
                    """
                    INSERT INTO Lactation (Number, Start, "End", id_Cow, id_Lactation_phase)
                    VALUES (%s, %s, %s, %s, %s)
                    """,
                    (
                        lact_num,
                        start_date,
                        end_date,
                        cow_id,
                        phase_id
                    )
                )

                total_records['Lactation'] += 1

                # следующая фаза начинается после предыдущей
                current_date = end_date

            # небольшой перерыв между лактациями
            current_date += timedelta(days=random.randint(30, 90))

    conn.commit()
    print(f"Добавлено {total_records['Lactation']} записей лактации")


def create_feeding_plan():
    cursor.execute("SELECT id_Diet FROM Diet")
    diet_ids = [row[0] for row in cursor.fetchall()]

    cursor.execute("SELECT id_Feeds FROM Feeds")
    feed_ids = [row[0] for row in cursor.fetchall()]

    cursor.execute("SELECT id_Day_times FROM Day_Times")
    day_time_ids = [row[0] for row in cursor.fetchall()]

    if not diet_ids or not feed_ids or not day_time_ids:
        print("Пустые зависимости!")
        return

    for diet_id in diet_ids:

        for day_time in day_time_ids:

            meals_count = random.randint(2, 3)

            for _ in range(meals_count):

                cursor.execute(
                    """
                    INSERT INTO Feeding_plan (Feed_amount, id_Feeds, id_Diet, id_Day_times)
                    VALUES (%s, %s, %s, %s)
                    """,
                    (
                        round(random.uniform(0.5, 12.0), 3),
                        random.choice(feed_ids),
                        diet_id,
                        day_time
                    )
                )

                total_records['Feeding_plan'] += 1

    conn.commit()
    print(f"Добавлено {total_records['Feeding_plan']} планов кормления")

def create_nutrition():
    # Получаем ID всех коров
    cursor.execute("SELECT id_Cow FROM Cow")
    cow_ids = [row[0] for row in cursor.fetchall()]

    # Получаем ID всех существующих планов кормления
    cursor.execute("SELECT id_Feeding_plan FROM Feeding_plan")
    plan_ids = [row[0] for row in cursor.fetchall()]

    if not cow_ids or not plan_ids:
        print("Нет коров или планов кормления в базе данных!")
        return

    for cow_id in cow_ids:
        # Задаем случайное количество планов кормления (рационов) для текущей коровы от 20 до 30
        count_plans_for_cow = random.randint(20, 30)
        
        # Выбираем уникальные планы кормления для коровы.
        # min() защищает код от ошибки, если вдруг всего планов в базе будет меньше 30
        chosen_plans = random.sample(plan_ids, min(count_plans_for_cow, len(plan_ids)))

        for plan_id in chosen_plans:
            cursor.execute(
                """
                INSERT INTO Nutrition (id_Cow, id_Feeding_plan)
                VALUES (%s, %s)
                """,
                (cow_id, plan_id)
            )
            total_records['Nutrition'] += 1

    conn.commit()
    print(f"Добавлено {total_records['Nutrition']} записей в таблицу Питание")

create_cows() # коровы
create_feeds() # корма
create_diets() # рационы
create_lactations() # лактации
create_feeding_plan() # план кормления
create_nutrition() # питание
\end{Verbatim}